\begin{abstract}
\noindent
We present a Haskell implementation of the parametric lens framework for
gradient-based machine learning proposed by Cruttwell et al.
In this framework every component of a training pipeline---layer, loss
function, and optimiser---is a \emph{parametric lens}: a composable unit
carrying a forward pass and a manually derived backward pass as a single
algebraic structure.  Composition of lenses corresponds directly to the
wiring of modules in a computation graph, and the combined parameter type
is inferred automatically by the type system.

The implementation extends the original Python proof-of-concept in several
directions.  Tensor shapes, mini-batch sizes, compute devices, and element
dtypes are encoded as GHC type-level parameters, so architectural
mismatches are compile-time errors rather than runtime failures; all models
are written in batched form from the ground up, whereas the original
processes a single example at a time.  Autoencoders are expressed as two
ordinary lens pipelines composed end-to-end, demonstrating that the
framework is not limited to discriminative tasks.  Residual skip
connections are realised as a higher-order combinator \texttt{skipPara},
built entirely from existing lens primitives, with the identity gradient
path emerging automatically from the combinator structure rather than from
a separate derivation.  Scaled dot-product self-attention and its
multi-head generalisation are implemented as \texttt{ParaLens'} values with
fully hand-derived backward passes, including the rank-one Jacobian
correction for the softmax non-linearity, with the head-count and
embedding-dimension compatibility constraint enforced at the type level.
Variable-depth residual networks are constructed by typeclass induction
over type-level Peano naturals.  Zero overhead is verified by GHC Core
inspection: the entire lens abstraction---composition,
reparametrisation, the van Laarhoven \texttt{forall}---is absent from the
optimised output.
\end{abstract}
